\documentclass[11pt,a4paper]{article}
\usepackage[italian]{babel}
\usepackage[margin=2.5cm]{geometry}
\usepackage{parskip}
\usepackage{amsmath, amssymb, amsthm}
\usepackage[most]{tcolorbox}
\usepackage{array}
\usepackage{tikz}
\usetikzlibrary{decorations.pathreplacing}
\usepackage{hyperref}

% --- Riquadri con bordo nero, senza numero ---
% Uso: \begin{definizione} ... \end{definizione}
%      \begin{teorema}[Nome] ... \end{teorema}  (il nome è facoltativo)
\tcbset{nota/.style={colback=white, colframe=black, boxrule=0.5pt,
  sharp corners}}
\NewTColorBox{definizione}{}{nota,
  before upper={\textbf{Definizione.}\ }}
\NewTColorBox{teorema}{o}{nota,
  before upper={\textbf{Teorema\IfValueT{#1}{ (#1)}.}\ }}
\NewTColorBox{proposizione}{o}{nota,
  before upper={\textbf{Proposizione\IfValueT{#1}{ (#1)}.}\ }}
\NewTColorBox{assioma}{o}{nota, boxrule=1pt,
  before upper={\textbf{Assioma\IfValueT{#1}{ (#1)}.}\ }}

% --- Ambienti semplici (senza riquadro) ---
\theoremstyle{definition}
\newtheorem*{esempio}{Esempio}
\newtheorem*{osservazione}{Osservazione}
\newtheorem*{esercizio}{Esercizio}

% V e F nelle tavole di verità
\newcommand{\V}{\textsf{V}}
\newcommand{\F}{\textsf{F}}

\title{Estremo superiore ed estremo inferiore}
\author{Marco Cerbella}
\date{Lezione 2 -- 28 settembre 2026}

\begin{document}
\maketitle

\section{Estremo superiore ed estremo inferiore}

In $[0, 1)$ i maggioranti sono $[1, +\infty)$, e il più piccolo di essi è
$1$: è proprio il valore a cui l'insieme ``si avvicina''. Questo suggerisce
la definizione seguente.

\begin{definizione}
Sia $A \subseteq \mathbb{R}$ non vuoto e limitato superiormente.
L'\emph{estremo superiore} di $A$, indicato con $\sup A$, è il
\textbf{minimo dei maggioranti} di $A$.

Analogamente, se $A$ è non vuoto e limitato inferiormente, l'\emph{estremo
inferiore} di $A$, indicato con $\inf A$, è il \textbf{massimo dei
minoranti} di $A$.
\end{definizione}

\begin{center}
\begin{tikzpicture}[scale=1.5]
  \draw[->] (-0.3,0) -- (7,0) node[right] {$\mathbb{R}$};
  \draw[line width=3pt, black!25] (0.5,0) -- (3.5,0);
  \fill (0.5,0) circle (1.8pt);
  \draw[fill=white] (3.5,0) circle (1.8pt);
  \draw[decorate, decoration={brace, amplitude=5pt}]
    (0.5,0.2) -- (3.5,0.2) node[midway, above=6pt] {$A$};
  \draw[line width=3pt, black!60] (3.5,-0.35) -- (6.8,-0.35);
  \node[below] at (5.2,-0.45) {maggioranti};
  \draw[thick] (3.5,0.15) -- (3.5,-0.6);
  \node[above] at (3.5,0.55) {$\sup A$};
\end{tikzpicture}
\end{center}

\begin{esempio}
\leavevmode
\begin{enumerate}
    \item $A = [0, 1)$: $\sup A = 1$ (che non è massimo, perché $1 \notin A$)
          e $\inf A = \min A = 0$.
    \item $A = \left\{ 1 - \frac{1}{n} \;\middle|\; n \in \mathbb{N},\
          n \geq 1 \right\} = \left\{ 0, \frac12, \frac23, \frac34, \dots
          \right\}$: gli elementi crescono avvicinandosi a $1$ senza mai
          raggiungerlo. Quindi $\sup A = 1$, ma $A$ non ha massimo;
          $\inf A = \min A = 0$.
\end{enumerate}
\end{esempio}

\textbf{Legame con massimo e minimo.}
\begin{itemize}
    \item Se $A$ ha massimo, allora $\sup A = \max A$.
    \item Se $\sup A \in A$, allora $\sup A$ è il massimo di $A$. Se invece
          $\sup A \notin A$, allora $A$ non ha massimo.
    \item Lo stesso vale per $\inf A$ e il minimo.
\end{itemize}

\begin{teorema}[Caratterizzazione dell'estremo superiore]
Sia $A \subseteq \mathbb{R}$ non vuoto e limitato superiormente. Allora
$s = \sup A$ se e solo se:
\begin{enumerate}
    \item $a \leq s$ per ogni $a \in A$ \quad ($s$ è un maggiorante);
    \item per ogni $\varepsilon > 0$ esiste $a \in A$ tale che
          $a > s - \varepsilon$ \quad (nessun numero più piccolo di $s$ è
          un maggiorante).
\end{enumerate}
\end{teorema}

In parole: $s$ sta sopra tutti gli elementi di $A$, ma se lo si abbassa
anche di pochissimo si trova subito un elemento di $A$ che lo supera. Per
l'estremo inferiore vale lo stesso con le disuguaglianze rovesciate:
$i \leq a$ per ogni $a \in A$ e, per ogni $\varepsilon > 0$, esiste
$a \in A$ con $a < i + \varepsilon$.

\section{Esistenza dell'estremo superiore}

Finora abbiamo dato per scontato che il minimo dei maggioranti esista. Non è
ovvio: l'insieme dei maggioranti potrebbe non avere minimo, come $[0, 1)$ non
ha massimo. È qui che interviene l'assioma di completezza.

\begin{teorema}[Esistenza dell'estremo superiore]
Ogni insieme $A \subseteq \mathbb{R}$ non vuoto e limitato superiormente
ammette estremo superiore in $\mathbb{R}$.

Analogamente, ogni insieme non vuoto e limitato inferiormente ammette
estremo inferiore in $\mathbb{R}$.
\end{teorema}

\begin{proof}
Sia $B$ l'insieme dei maggioranti di $A$. $B$ non è vuoto, perché $A$ è
limitato superiormente, e per definizione di maggiorante
\[
a \leq b \quad \text{per ogni } a \in A \text{ e per ogni } b \in B.
\]
Per l'assioma di completezza esiste un elemento separatore $c$:
\[
a \leq c \leq b \quad \text{per ogni } a \in A \text{ e per ogni } b \in B.
\]
Dalla prima disuguaglianza, $c$ è un maggiorante di $A$, cioè $c \in B$.
Dalla seconda, $c$ è minore o uguale a ogni maggiorante. Quindi $c$ è il
minimo di $B$, cioè $c = \sup A$.
\end{proof}

\textbf{In $\mathbb{Q}$ il teorema è falso.} L'insieme
$A = \{q \in \mathbb{Q} \mid q > 0,\ q^2 < 2\}$ è limitato superiormente
(ad esempio da $2$), ma tra i \emph{razionali} non ha estremo superiore:
l'unico candidato è $\sqrt{2}$, che non è razionale. Per questo l'esistenza
dell'estremo superiore è una forma equivalente dell'assioma di completezza.


\textbf{Convenzione.} Se $A$ non è limitato superiormente si scrive
$\sup A = +\infty$; se non è limitato inferiormente si scrive
$\inf A = -\infty$. Ad esempio $\sup \mathbb{N} = +\infty$ e
$\inf \mathbb{Z} = -\infty$.

\section{Esempi}

\begin{esempio}[L'insieme $B$]
$\max B = \sup B = 7$. Per l'estremo inferiore verifichiamo che $\inf B = 3$:
\begin{enumerate}
    \item $3 \leq x$ per ogni $x \in B$: $3$ è un minorante;
    \item per ogni $\varepsilon > 0$ esiste $x \in B$ con $x < 3 + \varepsilon$:
          infatti $[3, 3 + \varepsilon[ \,\cap\, ]3, 7] \neq \emptyset$.
\end{enumerate}
Il numero $3$ non appartiene a $B$, quindi $B$ non ha minimo. Invece $0$ è
un minorante ma \textbf{non} è l'estremo inferiore: per $\varepsilon = 2$ non
esiste nessun $x \in B$ con $x < 0 + 2$ (la condizione 2 fallisce).
\end{esempio}

\begin{esempio}[$A = \{1/n \mid n \in \mathbb{N} \setminus \{0\}\}$]
$A = \{1, \frac12, \frac13, \frac14, \dots\} \subseteq [0, 1]$, quindi $A$ è
limitato.
\begin{itemize}
    \item $\max A = \sup A = 1$: infatti $1 \in A$ (è $x_1$) e
          $x_n = \frac1n \leq 1$ per ogni $n \geq 1$.
    \item $\inf A = 0$. Si ha $0 \leq x_n$ per ogni $n$. Dimostriamo che per
          ogni $\varepsilon > 0$ esiste $x_n \in A$ con $x_n < 0 + \varepsilon$.
          Per assurdo, supponiamo che esista $\varepsilon > 0$ tale che
          $\frac1n \geq \varepsilon$ per ogni $n \in \mathbb{N} \setminus \{0\}$;
          allora $n \leq \frac1\varepsilon$ per ogni $n \in \mathbb{N}$, che è
          assurdo perché $\mathbb{N}$ non è limitato superiormente. Quindi
          $0 = \inf A$.
    \item $0 \notin A$, quindi $A$ \textbf{non ha minimo}.
\end{itemize}
\end{esempio}

\textbf{Riassunto.}
\begin{itemize}
    \item se un insieme ha massimo, allora massimo $=$ estremo superiore;
    \item se un insieme ha minimo, allora minimo $=$ estremo inferiore;
    \item un insieme può avere estremo inferiore ma non minimo;
    \item un insieme può avere estremo superiore ma non massimo.
\end{itemize}
\end{document}
